

\chapter{Introduction}

    The use of Field Programmable Gate Arrays (FPGAs) as real-time data
    processors is attracting growing interest. The
    Xilinx Virtex FPGA series is aimed directly at the DSP market
    and promises spectacular performance in comparison to a conventional
    CPU.

    However, the programming of FPGAs still presents considerable
    difficulties, and digital circuit design is still the predominant
    method. Parallel languages for FPGAs have emerged in recent years, but
    they have arisen largely from computer science departments where the
    focus has been on a wide range of general parallel algorithms. The
    environment of parallel CPUs which prompted early work, the very
    generality of these parallel languages, and their evolution from serial
    languages has resulted in compilers that arguably do not use FPGAs very
    efficiently. In particular these parallel programming languages do
    not support pipelining very well.

    Pipeline processing seems a natural way to implement machine vision
    (and other) algorithms which work on real-time data streams. Even
    where data is stored in memory, in many cases it can be traversed
    to generate a data stream. Even where an algorithm is sequential,
    pipelining results in some parallel execution - the parallelism
    is effectively the length of the pipeline.

    This software was developed in 2001 within a machine vision research
    group in CSIRO (Commonwealth Scientific and Industrial Research
    Organisation), an Australian Federal Government research laboratory.
    It was developed and used successfully over two decades to implement
    machine vision algorithms in FPGAs and to provide high-speed data
    processing in instrumentation systems.

    3PL has now been released as open-source software. As it was developed as
    an experimental research tool it carries no guarantees. Although it was
    refined over a long period, it will almost certainly not be error free and
    thus {\bf 3PL must not be used in any critical applications!}

    3PL is implemented in Java and so will operate under both UNIX/LINUX/OSX
    and Windows environments. The release includes the source code however
    the Java jar file is provided so it is unnecessary to rebuild 3PL unless
    a bug is to be corrected or some extension or improvement is contemplated.

    \section{The 3PL language}

        3PL is a {\em dataflow}\footnote{Dataflow was a name
        given to data-driven computing architectures explored at Manchester
        University by Ian Watson and John Gurd in the late 1970s.} parallel
        programming language for FPGAs which has been explicitly designed to
        support pipelining. The name {\bf 3PL} stands for "Parallel Pipeline
        Programming Language".

        In principle we would like our pipeline to have the following
        characteristics -
        \begin{enumerate}
        \item    Each pipeline module should process input data when -
                \begin{enumerate}
                \item    its internal state is appropriate
                \item    required input data are available
                \item    required output can take place
                \end{enumerate}
        \item   Movement of data through the pipeline should depend only
                on the behaviour of individual modules and their
                interaction.
        \end{enumerate}

        Data will flow through the pipeline as rapidly as possible, retarded
        only when one or more source module queues are empty or a destination
        module queue is full. Thus such pipelines do not necessarily move
        in lock-step with the clock.

        {\bf 3PL} supports pipelined data, and modules are inherently
        parallel. While the pipeline/module model works well for real-time
        data streams, there are some sequential algorithms for which pipeline
        processing is clumsy and laborious. While non-pipelined sequential
        operations are inefficient in an FPGA (a modest conventional CPU will
        be significantly faster), occasions arise where the sequential code is
        but a small part of a larger, predominantly parallel, algorithm. For
        this reason {\bf 3PL} supports the usual conventional sequential
        code blocks as well as parallel code blocks.

        Sequential constructs can also be used in code which is executed at
        compile time, rather than in the target architecture. This {\em
        immediate} code is used to control target code generation, for example
        to generate a network topology from an algorithm or to systematically
        initialise a table.

        The 3PL compiler is not specific to any particular FPGA, but the code
        generator supports only Xilinx FPGAs and CPLDs at the moment.
        Currently these platforms are Spartan2, Spartan3, Spartan6, Virtex,
        Virtex-E, Virtex2, Virtex2-Pro, Virtex4, Virtex5, Virtex6 and series 7
        FPGAs (Artix, Kintex, Virtex7 and Zynq). It also supports Coolrunner
        and XC95 CPLDs though only a much restricted range of 3PL features can
        be used for these devices.

    \section{Current version}

        The current version of 3PL is {\bf 11.7}. It requires Java 1.6 or
        later. A list of compiler versions appears as an appendix.

    \section{Known deficiencies}

        \subsection{enumerated types where ordinal 0 is not used}

            When a queue variable whose type is non-primitive is assigned, any
            words in the type that are not assigned in that clock cycle are by
            default assigned a value corresponding to all bits low. For numeric
            types that is zero and for logical that is false. This is a useful
            convention that can be employed to simplify source code. For an enum
            that default will be the value which has ordinal 0. If an enum is
            created using inbuilt procedure enumv() then ordinal 0 may not be
            used. In that case the above scenario would result in an invalid
            enum value. This problem can be easily avoided by either using
            ordinal 0 or else always explicitly assigning that field, so it
            is unlikely that this will be fixed.

        \subsection{++ and -- operators limited}

            The increment and decrement operators work on integer primitives,
            including members of arrays and structs. They do not parse if
            applied to an expression and hence cannot be used on an indirect
            reference via the \verb'->' operator.

        \subsection{cannot have arrays of queues, clocks, memories etc}

            Arrays and structs are part of a type specification. Clock mode
            variables do not have a type and hence there cannot be an array of
            clocks. The two memory modes do not have types, though their address
            and data types must be supplied, and so again there cannot be an
            array of these. Priority mode variables have a type which is "[]log"
            but again there cannot be an array of priority variables. Queue mode
            variables can have an array type, but the entire variable is read
            and written as a unit. You can have a queue whose type is an array
            but you cannot have an array of queues.

            Immediate, value, selectvalue and static mode variables can be
            assigned or evaluated in part or whole and hence do not present a
            problem in this respect.

            This inconvenience is a result of the language semantics and is not
            simply an implementation deficiency.

            Arrays of any mode of variable can be created and accessed by means
            of pointers. An array of pointers is created and then initialised to
            point to created variables. This can be done in a function so that
            the created variables are hidden and the pointer array is returned
            as the value of the function. There is an example in library
            stdlib.3pl, {\bf queuearray()}, which creates a number of queue mode
            variables and returns an array of pointers to the queues.

        \subsection{petrinets transitions do not wait on code blocks}

            When a petrinet transition occurs an associated code block may wait
            on queue availability. Transitions should wait for completion of
            associated code blocks. To implement this the transition would have
            to occur simultaneously with the last execution cycle of the code
            block. At present there is no easy way to do this. This will be
            further considered since this problem limits the applicability of
            petrinets.






    \section{Future development}

        \subsection{static multiply-assignment operator}

            The \verb'*<-' assignment may be implemented later.

        \subsection{extension to round function}

            The round function takes both immediate and target mode arguments.
            For the latter it rounds a fixed to int or ufixed to uint. It may
            be extended in the future to allow rounding of fixed, ufixed, int
            or uint at any bit position, producing an appropriate result type,
            but so far no particular need for this has been felt.

        \subsection{target code pointers}

            Immediate-mode pointers are available in 3PL. In immediate code,
            pointers to immediate or target mode variables can be created,
            assigned and dereferenced. There is no equivalent in target code.
            Pointer variables which could be assigned and dereferenced in target
            code could be implemented, but the resulting decoder logic would be
            extensive and slow except in the case of very limited usage, hence
            this is probably not worth considering.

        \subsection{module and procedure call syntax}

            Currently the syntax for calling a module is identical to that for a
            procedure. It would enhance readability if these calls had
            distinctive syntaxes. This would help emphasise the difference
            between a module and a procedure. In particular the fact that a
            called module is a completely separate (new) environment from that
            of the call. I am considering changing module calls to the following
            syntax -
            \begin{lstlisting}[gobble=12, language=threepl]
               [modname(arg, arg, ....)]
            \end{lstlisting}
            using the brackets to suggest a box as for inline modules.
            Note that an inline module containing a procedure call is
            distinguished by the ";" terminating the call -
            \begin{lstlisting}[gobble=12, language=threepl]
               [procname(arg, arg, ....);]
            \end{lstlisting}
            The proposed change has not met with enthusiasm from other users
            so it may never be done.

        \subsection{debugging facilities}

            Six debugging tools have been envisaged -
            \blist
            \item   an immediate execution interactive debugger
            \item   a target code simulator
            \item   a timing constraint error report analyser which will
                    determine which parts of the source are responsible for
                    failure to meet place-and-route timing constraints
            \item   a target run-time exception mechanism which can report
                    exceptions or freeze execution when they occur
            \item   a single step facility, at least for single clock
                    applications
            \item   a configuration read-back utility which would relate
                    variable and memory contents back to the source code
            \blend

            The second, a target simulator, has been included as part of the
            compiler/interpreter. As 3PL evolved the simulator atrophied
            and hence it has been disabled, though the source code is still
            included to allow later re-enabling if required.

            Provision of the other five facilities has not yet been started. The
            third tool listed is the most urgent because failure to meet timing
            constraints will be a major problem for users of the language, but
            it has not yet been tackled.

        \subsection{variable types}

            Thus far bit, integer, unsigned integer, fixed point, floating point
            (limited), enumerated and logical types have been implemented for
            target variables. The string type can be used for immediate
            variables but there is no plan to extend it to target variables.

        \subsection{linkage to precompiled 3PL code}

            No mechanism for this has been designed or implemented.
            Currently all 3PL libraries are in 3PL source code and are
            compiled along with the application program.



    \section{Acknowledgements}

        This work was done within what was the Machine Vision group in CSIRO
        \footnote{Commonwealth Scientific and Industrial Research Organisation -
        an Australian federal government research laboratory}
        Manufacturing and Infrastructure Technology, Preston laboratory,
        Melbourne. The group progressed to become known as High Speed
        Instrumentation, moving to CSIRO Mineral Resources, Clayton, Melbourne.
        The group evolved further and is now called MicroAnalytical
        GeoScience (MaGs).

        Paul Dunn conceived the language and he and Andrew Bainbridge-Smith
        developed the original (incomplete) compiler. Andrew did all the
        groundwork finding java tools and set up the make file, the directory
        structure and the javadoc framework and wrote the original parser. Paul
        wrote the code generator.

        The first version of the compiler was a good learning exercise but soon
        proved to have serious shortcomings and was never completed. Andrew left
        the laboratory shortly thereafter and Paul redesigned the language and
        re-implemented the compiler. This benefited greatly from the experience
        of the first attempt.

        Andrew Tulloh has worked through both compiler versions developing the
        target simulator (currently disabled as it has not kept up with 3PL
        evolution).

        Robin Kirkham has been an invaluable critic and expert user of 3PL
        throughout and has contributed significant 3PL library code.

        Murray Jensen crafted the make files to operate under Unix-like
        or Windows environments (development and subsequent usage of 3PL
        has however been entirely under UNIX/LINUX/OSX). In later years,
        as the compiler is predominantly a Java program, Murray switched
        the build system over to Apache Ant, which seemed to offer more
        Java oriented dependency support. It was hoped that this would
        handle things like inner-classes and jar files better (it is
        unclear whether this has been successful).
